\section{灵敏度分析与误差分析}

为检验模型结果对关键参数与输入误差的依赖程度，本章从灵敏度分析与误差分析两个方面进行检验：
灵敏度分析考察关键参数变化时最优方案与购电费用的变化幅度，用以判断模型结论是否稳定；
误差分析考察预测值与实际值之间的偏差及其经济后果，用以判断结果是否可信。

\subsection{灵敏度分析}

\subsubsection{预测方法与余量参数}
模型中不确定性最强、对结果影响最大的参数是预测方法与预测余量系数。
表~\ref{tab:sens-q2} 给出问题二在不同预测方法与余量组合下的全年总费用与紧急购电费用。

\begin{table}[!htbp]
\caption{不同预测方法与余量组合下的全年运行费用（问题二，报告期 334 天）}\label{tab:sens-q2}
\centering
\small
\setlength{\tabcolsep}{5pt}
\begin{tabular}{llcc}
\toprule
预测方法 & $(\gamma,\delta)$ & 总费用/万元 & 紧急购电费/万元 \\
\midrule
持续性预测（当日 $=$ 昨日） & $(1.00,1.00)$ & 2244.8 & 1021.1 \\
持续性预测 & $(1.05,1.00)$ & 2139.4 & 793.6 \\
上周同星期（w1） & $(1.00,1.00)$ & 1602.4 & 378.7 \\
\textbf{上周同星期（w1，本文采用）} & $\bm{(1.05,0.98)}$ & \textbf{1454.7} & \textbf{87.0} \\
近 4 个同星期均值 & $(1.05,1.00)$ & 1574.7 & 229.7 \\
7 日滑动平均 & $(1.08,1.00)$ & 2059.0 & 698.6 \\
\bottomrule
\end{tabular}
\end{table}

由表~\ref{tab:sens-q2} 可知，预测方法的选取对费用的影响远大于余量系数：
持续性预测会因周末前后的负载状态切换而严重失效，全年总费用高达 2244.8 万元；
若完全不加安全余量，w1 预测的费用为 1602.4 万元，加入余量后降至 1454.7 万元，
降幅达 147.7 万元，说明\textbf{余量策略是应对 5 倍紧急购电惩罚的关键}。
在 w1 预测下，$\gamma$ 在 $[1.03,1.07]$ 区间内构成平坦平台，即 $\gamma$ 偏移 $\pm0.02$ 时总费用变化不足 1\%；
将报告期拆分为 2--8 月与 9--12 月两段分别寻优（时间序列交叉验证的思想\cite{fpp3} 第 5.6 节），得到的最优参数与全年结果一致，
说明该参数并非针对全年数据"调"出来的，模型未发生过拟合。

从理论上看，5 倍惩罚使购电决策具有强不对称性：多购 1\,kWh 最多浪费约 1 倍电价，
少购 1\,kWh 则需按 5 倍电价结算，这正是\textbf{报童模型}的结构，
其临界分位数为 $(5-1)/5=0.8$，即日前计划应覆盖净负荷约 80\% 分位数而非中位数。
按全年均值折算，取 $\gamma=1.05$、$\delta=0.98$ 时对净负荷的等效余量约为 12\%，
与临界分位数所要求的安全水平量级一致，说明该参数具有明确的理论依据。

\subsubsection{策略参数的稳健性}
\textbf{（1）储能配置的价值。} 在同一预测与余量下取消储能设备，问题二全年总费用由 1454.7 万元升至 1893.9 万元，
即储能带来的年价值为 439.2 万元（降低 23.2\%），说明储能是实现电能在时间上转移、消化预测偏差的核心手段。

\textbf{（2）当日残差更新权重。} 在问题三中，权重取 $\lambda=0$（不使用当日已实现信息）时全年总费用为 1376.3 万元，
取最优值 $(\lambda_{pv},\lambda_L)=(0.1,0.3)$ 时降至 1368.8 万元，
而权重继续增大至 $(0.6,0.8)$ 时费用回升至 1414.4 万元（见表~\ref{tab:q3-lambda}）。
这说明当日信息确有价值，但权重过大会使调整决策追逐残差中的随机噪声、支付多余的调整费用；
最优解附近费用变化平缓，说明结果对 $\lambda$ 的具体取值不敏感。

\textbf{（3）末端最低备用约束。} 该约束在 2000$\sim$7000\,kWh 范围内对应的全年总费用变化平缓，
即备用水平的选取不改变"标定 $+$ 双向调整优于无调整"这一结论；
本文取 4000\,kWh（约为可用容量的 42\%）主要出于稳健性考量——避免储能被放空而丧失次日备用能力。

\textbf{（4）波动电价下的参数复核。} 引入逐日波动电价（问题四）后重新扫描：
在问题二策略下，$\gamma$ 取 1.03、1.05、1.07 时全年总费用分别为 1530.2、1526.4、1548.6 万元，
最优值仍在 $\gamma=1.05$ 附近，故该策略继续沿用 $(\gamma,\delta)=(1.05,0.98)$；
在问题三策略下，负载侧最优值仍为 $\gamma=1.02$，而光伏侧余量由 $\delta=0.98$ 调整为 $\delta=1.00$，
对应 $\delta$ 取 0.96、0.98、1.00 时费用分别为 1440.3、1433.6、1431.6 万元。
这说明电价波动改变了最优光伏余量（波动电价下调价空间更大，过度下压光伏预测反而抬高了购电成本），
但并未改变模型结构与参数的整体结论。

\subsection{误差分析}

\subsubsection{光伏预报误差}
对附件 3 的整点预报与实际出力的对照（图~\ref{fig:q3-forecast}）表明，
预报平均绝对误差在 108.0\,kW 至 1374.1\,kW 之间，并呈现稳定的\textbf{系统性偏差}：
上午时段预报偏低（8:00 约 $-1050$\,kW），午后时段预报偏高
（15:00 至 16:00 约 $+836$\,kW 至 $+1067$\,kW），且预报峰值时刻与实际峰值时刻存在错位。
经 MOS 标定后，14:00 的估计 MAE 由原始预报的 793.4\,kW 降至 373.0\,kW，
也优于单独使用上周同星期实际值的 456.8\,kW，说明\textbf{标定有效削弱了系统性偏差}。
由于预报误差不可能完全消除，本文并未追求更高的预测精度，而是通过安全余量与日内滚动调整来
吸收剩余误差，这与误差分析的目的（判断结果是否可信）是一致的。

\subsubsection{预测误差的经济代价}
为量化预测误差带来的经济损失，本文以"完美预测"作为理论下界：
在与最终策略完全相同的框架下（含储能、日循环约束与跨日滚动），
仅将预测值替换为当日实际值、并将余量置为 1，重新求解全年调度。
结果表明，完美预测下问题二的全年总费用为 1224.5 万元，
而本文最终策略为 1454.7 万元，二者相差 230.2 万元，即\textbf{残余预测误差的代价约为总费用的 15.8\%}。
在问题三中，由于引入预报标定与日内滚动调整，该差距缩小至 144.3 万元（1368.8 万元对 1224.5 万元），
说明滚动调整机制确实降低了预测误差造成的经济损失。

进一步考察误差代价的分布可以发现，问题二中紧急购电的平均结算电价为 1.094 元/kWh，
约为全天平均电价的 1.4 倍，说明供电缺口在时间上并非均匀分布，
而是较多落在傍晚高电价时段——5 倍电价惩罚在最贵的时段被放大，
这也是问题三引入日内滚动调整所能获得收益的主要来源。

\subsubsection{余量策略的固有代价与数值校验}
日前计划在预测值上施加安全余量，其代价是部分已付费电量最终无法使用：
问题二全年弃电量约 326 万\,kWh（占计划购电量的 14.6\%），
其中"已付费却被迫丢弃"的电量下限约 81 万\,kWh，按均价估算浪费不低于 60 万元（约占总费用 4\%）。
这一弃电是“日前承诺 $+$ 不确定性对冲”的必然成本，
但与所避免的约 790 万元紧急购电费用相比，净收益仍然显著。

最后，本文对求解结果进行全链路数值校验，主要包括以下五个方面：
\begin{enumerate}[label=(\arabic*), leftmargin=2.5em, itemsep=0.3em]
\item 各时段功率平衡关系与储能状态递推关系均成立；
\item 储能轨迹全程位于 $[1200,10800]$\,kWh 区间内，充放电功率不超过 5000\,kW；
\item 跨日储电量按实际末态连续传递，无跳变；
\item 各日费用按“计划购电费 $+$ 调整费用 $+$ 紧急购电费”分解后与总费用一致；
\item 调整模型的解可由重新求解线性规划复现。
\end{enumerate}
上述校验在 334 天的滚动仿真中均未发现违规，说明模型的求解结果在数值上是自洽且可复现的。

